% Rozklad tyče na první díl délky $i$ a zbytek délky $n-i$. Zbytek je
% opět úloha téhož druhu, jen menší.
\begin{tikzpicture}[x=1cm,y=1cm,font=\small,
                    dil/.style={thick, fill=lightred!25, rounded corners=1pt},
                    zbytek/.style={thick, fill=darkgreen!20, rounded corners=1pt}]

    % Tyč délky n rozdělená řezem
    \draw[dil]    (0,0)   rectangle (2.4,0.55);
    \draw[zbytek] (2.4,0) rectangle (8.0,0.55);

    \node at (1.2,0.275) {$i$};
    \node at (5.2,0.275) {$n-i$};

    % Řez
    \draw[darkred, line width=1pt, dashed] (2.4,-0.30) -- (2.4,0.85);
    \node[darkred, font=\scriptsize, anchor=south] at (2.4,0.90) {řez};

    % Popisky
    \draw[<->,>=stealth] (0,-0.45) -- (2.4,-0.45);
    \draw[<->,>=stealth] (2.4,-0.45) -- (8.0,-0.45);

    \node[font=\scriptsize, anchor=north, align=center, text width=3.0cm,
          lightred!70!black] at (1.2,-0.60)
        {první díl prodáme\\ za cenu $p_i$};
    \node[font=\scriptsize, anchor=north, align=center, text width=5.0cm,
          darkgreen!70!black] at (5.2,-0.60)
        {zbytek je táž úloha o~velikosti $n-i$\\ $\Rightarrow$ nejlépe za $r_{n-i}$};

    \node[font=\small, anchor=north] at (4.0,-1.85)
        {$r_n=\max\limits_{1\leqslant i\leqslant n}\left(p_i+r_{n-i}\right)$};

\end{tikzpicture}
